% =============================================================
% Conclusion générale (~2 pages) — sans numérotation de section.
% Cinq points imposés (note pédagogique ESPRIT) : récapitulation,
% résultats chiffrés, problèmes rencontrés, apports (entreprise et
% élève-ingénieur), perspectives.
% Aucun chiffre de cette page n'est nouveau : tous sont repris du
% corps, avec leur date.
% =============================================================
\chapter*{Conclusion générale}
\addcontentsline{toc}{chapter}{Conclusion générale}
\markboth{Conclusion générale}{Conclusion générale}

\textbf{Récapitulation.} Ce mémoire a présenté la conception, la réalisation et la validation
d'un socle d'hébergement cloud Zero Trust pour une entreprise émergente. Le point de départ
était une application fonctionnelle, mais hébergée sur un serveur unique, avec confiance
implicite entre tous ses composants, sans reproductibilité et sans détection de sécurité. Le
point d'arrivée est un socle entièrement décrit en code, organisé en sept niveaux L1--L7, soit
cinq couches et deux plans transversaux. Il comporte deux chaînes de livraison à portes
bloquantes, une chaîne de détection enrichie par un modèle de représentation sémantique, et un
tableau de bord destiné à l'analyste.

\textbf{Résultats chiffrés.} Le socle fonctionne avec deux environnements réels, développement
et recette ; la production reste une cible non provisionnée, et l'infrastructure durable du
périmètre est décrite en code avec ses décisions et ses écarts datés. La restauration de la base
de données, mesurée le 03/08/2026 sur la configuration zonale antérieure au passage en haute
disponibilité régionale du 08/08/2026, a demandé 32~min~45~s pour 314~lignes restaurées sur 314
au point demandé, sans perte de données. Deux retours arrière applicatifs ont pris 11,6~s et
16,5~s le 07/08/2026 ; le démarrage à froid du service d'encodage est de 27~s en moyenne et de
94~s au pic. Le 19/08/2026, un scénario de bout en bout a produit treize~requêtes malveillantes,
toutes bloquées, agrégées en une alerte environ quinze~minutes plus tard. La campagne de
validation établit \textbf{sept résultats conformes, onze partiellement conformes et deux
non-conformités assumées} : les vingt protocoles portent une preuve nommée, et six d'entre eux
une cadence de rejeu au plan de validation continue, dont le premier relevé est daté.

Trois de ces relevés sont postérieurs à la migration du réseau --- un refus de sortie réellement
opposé et journalisé, un refus croisé entre locataires, et la non-conformité du composant
d'inférence, dont la sortie échappe encore au pare-feu. Le compte est volontairement sévère :
n'est déclaré
conforme qu'un critère couvert intégralement, le plus souvent par un refus réellement observé.
La couche sémantique, elle, produit des candidats de rattachement au positionnement
informationnel : sa précision comparative et l'effet de la boucle F6 sur l'ordre de priorité
relèvent du même plan. La seconde moitié de la problématique attend donc encore sa mesure --- le
prix payé en faux positifs et en latence de détection demande une période d'observation sous
trafic réel, et c'est la première des perspectives.

\textbf{Problèmes rencontrés.} \textbf{Cinq} incidents d'ingénierie ont été documentés au
\cref{ch:realisation}. Le premier est une indisponibilité liée au démarrage à froid. Le deuxième est un
blocage collectif d'utilisateurs légitimes par la limitation de débit. Le troisième est un état
local divergent après une application partiellement échouée. Le quatrième est une porte
d'analyse statique restée verte sans évaluer aucune règle. Le cinquième est une règle de sortie
qui n'aurait rien gouverné : inscrite au plan de correction, elle eût été sans effet, le trafic
ne passant pas là où elle se serait appliquée --- d'où l'ordre de la migration réseau, déplacer
le trafic, l'observer sept~jours, n'interdire qu'ensuite. Chacun a produit une correction et une
leçon transférable ; la dernière est la plus générale, \textbf{un contrôle qui ne voit pas passer
le trafic n'est pas un contrôle, c'est une déclaration}.

Deux limites sont énoncées avec leur condition de levée. Le contrôle dynamique des sorties réseau
a été mesuré les 24 et 25/08/2026 :
il refuse ce qu'il gouverne, et deux charges de travail détachées du sous-réseau lui échappent
encore par une sortie managée --- non-conformité T14 du 25/08/2026. Surtout, l'enrichissement
sémantique occupe un positionnement informationnel assumé : il n'alimente pas le score
d'incident, et sa condition d'intégration est nommée --- une clé d'entité dans le schéma
d'enrichissement, une règle de décision, et la mesure comparative qui l'autorise.

\textbf{Apports.} Pour l'entreprise, le socle fournit une base interne réutilisable, éprouvée
partiellement par l'accueil d'un second locataire le 07/08/2026. L'identité, les secrets et les
données --- et désormais le réseau, le refus croisé ayant été réellement opposé le 24/08/2026 ---
sont séparés par locataire ; l'instance de base de données et le jeu de tables de supervision,
eux, restent partagés. Pour ma part, ce projet a mobilisé l'ensemble des compétences d'un
ingénieur en cybersécurité cloud, sans répartition d'équipe : audit, dérivation des exigences,
conception de l'architecture, écriture de l'infrastructure en code, construction des deux chaînes
de livraison et de la chaîne de détection, puis conduite de la campagne de validation. J'y ai
appris à documenter un écart au moment où il apparaît plutôt qu'à le refermer discrètement, et
à préférer une mesure défavorable à une estimation flatteuse --- le démarrage à froid mesuré à
27~s contre une estimation initiale de quelques secondes, et la porte d'analyse statique
découverte en faux vert, en sont les deux exemples les plus formateurs. J'y ai appris enfin qu'un
contrôle correctement conçu peut produire un effet non prévu en exploitation, comme cette
limitation de débit qui bloquait des utilisateurs légitimes partageant une même adresse.

\textbf{Perspectives.} Quatre axes d'amélioration sont identifiés. Le premier est
l'\textbf{isolation des données par locataire}, conçue, chiffrée et ordonnancée au \cref{ch:conception}
sous la décision D16 : ce qui reste est l'application, non l'instruction --- et la raison de
l'ordonnancement est écrite, deux migrations d'infrastructure dans la même quinzaine n'étant pas
soutenables par un opérateur unique. La segmentation réseau qui l'accompagnait, elle, a été
\textbf{appliquée du 11 au 25/08/2026} en douze étapes, et mesurée : le protocole T20 est
conforme au 24/08/2026 sur un refus croisé entre locataires réellement opposé et journalisé. Le
deuxième axe est l'intégration contrôlée de l'enrichissement au score d'incident, avec clé
d'entité, règle de décision et mesure comparative. Le troisième est la signature des images par
chaîne de provenance. Le quatrième est l'écriture de tests unitaires par règle de détection,
vérifiant que chaque règle produit l'alerte attendue sur des données de test. Le gain de
l'enrichissement sémantique, une fois mesuré, justifiera-t-il son coût ? C'est la question
ouverte du \cref{ch:etat-art}, et elle reste ouverte.

Ce travail ne produit pas un service commercialisable en l'état : il produit un socle
reproductible et mesuré, dont les limites sont chiffrées et dont le chemin vers un engagement
contractuel est identifié et estimé. C'est précisément parce que ces limites sont écrites que ce
chemin est crédible.
